\section{Deep RL 的两大算法范式}

在介绍了 RL 的基本概念之后，本节重点讨论如何用\textbf{深度神经网络}来实现强化学习算法。Deep RL 有两大主要范式：

\begin{figure}[H]
\centering
\includegraphics[width=0.85\textwidth]{slides-images/slide-20.jpg}
\caption{Deep RL 的两大算法范式：Value Learning 与 Policy Learning\protect\footnotemark}
\end{figure}
\footnotetext{Slides 第20页。}

\begin{importantbox}{Value Learning vs. Policy Learning}
\begin{itemize}
    \item \textbf{Value Learning}（价值学习）：先学习 Q-function $Q(s, a)$，再通过 $a = \arg\max_a Q(s,a)$ 确定最优动作
    \item \textbf{Policy Learning}（策略学习）：直接学习策略函数 $\pi(s)$，从策略中采样动作 $a \sim \pi(s)$
\end{itemize}
两种方法各有优劣，适用于不同类型的问题。
\end{importantbox}

\subsection{本章小结}
Deep RL 的核心思想是用深度神经网络作为函数逼近器，来学习 Q-function 或 Policy。Value Learning 通过间接方式（先学 Q 再决策），Policy Learning 通过直接方式（直接学策略）来解决决策问题。接下来两节将分别详细介绍这两种范式的代表性算法。


